\section{The Gaussian mean away from seeds}
\label{sec:local-mean}
The local covariance is accompanied by a local mean. Its defining
formula alone does not make the affine energy small: a uniform bound
on the inverse would only preserve the size of the original source.
The additional gain comes from the stationary linear response. We
first prove its cancellation, and then compare it with the actual
mean on one common clean neighbourhood.

\subsection{The stationary response cancels every linear row}
Fix the step, forest and full flat kernels of
Section~\ref{sec:seed-unhit-comparison}. Write $\mathcal U_F b$ for
the fine stationary field in tree coordinates with coarse input $b$.
It is the tree-straightened installed linear response. Its relative
field is
\[
 v_F(b)=\mathcal U_F b-\overline b,
\]
where $\overline b$ is the linear coarse replica. On every forest
bond this difference is zero. We use the actual diagonalization
identity for this response: if $K_f,K_c$ are the fine and coarse
quadratic forms and $\mathcal Q$ is the rooted linear predictor,
\begin{equation}\label{eq:mean-diagonal}
 K_c=\mathcal U_F^*K_f\mathcal U_F
       +p(1-\mathcal Q\mathcal U_F)^*
                     (1-\mathcal Q\mathcal U_F).
\end{equation}
These are the full physical forms, including the prescribed root
correction. This is an incoming flat construction identity; it does
not follow merely from positivity of the forms.

Let $b_s$ and $a_s$ be the pullbacks, by one signed lattice symmetry,
of $b$ and $\mathcal U_Fb$. Set
\[
 w=B_cd_cb_s,\qquad u=B_fd_fa_s,\qquad
 z=\sqrt p\,(b_s-\mathcal Qa_s),\qquad
 U=F_0w,\quad V=I_0w.
\]
All norms in this paragraph are counting Euclidean norms, summed
over their physical sites and colours. The two curvature squares
are the corresponding quadratic forms. Indeed expand each colour
component in the curvature factorization $K=d^*B^*Bd$; the predictor
square is exactly the last term in \eqref{eq:mean-diagonal}.
Thus that identity gives
\begin{equation}\label{eq:mean-norm-equality}
 \|u\|_2^2+\|z\|_2^2=\|w\|_2^2.
\end{equation}
The same equality holds after the signed pullback: the two full
quadratic forms are invariant, and the nonlinear predictor energy
is invariant with the transported tags. Apply that last invariance
to $\exp(tb)$ and $\exp(t\mathcal U_Fb)$, divide by $t^2$ and
let $t\to0$. The derivative proved in
Section~\ref{sec:predictor-branch} gives the stated predictor square.
This also explains why the rooted response need not itself commute
with the chosen symmetry.

The full stack identities give
\[
 N_0w=0,\qquad F_0^*U+I_0^*V=w,
 \qquad F_0^*u+I_0^*z=w.
\]
The first is the extra identity $N_0B_cd_c=0$; the second follows
from the stack isometry and the first, and the third is
\eqref{eq:unhit-linear-zero}. Taking inner products with $w$ gives
\[
 \|U\|_2^2+\|V\|_2^2=\|w\|_2^2,
 \qquad \langle U,u\rangle+\langle V,z\rangle=\|w\|_2^2.
\]
Together with \eqref{eq:mean-norm-equality}, these imply
$\|U-u\|_2^2+\|V-z\|_2^2=0$. Every summand is nonnegative,
so $U=u$ and $V=z$ at every address. In the notation of affine
energy rows, this says
\begin{equation}\label{eq:mean-stationary-zero}
 C_0(\overline{b_s},b_s)-G_0(b_s)
                 +\mathcal A_\infty(v_F(b))_s=0.
\end{equation}
Here $C_0$ is the linear fine-curvature/predictor pair, $G_0$ the
linear guide, and $\mathcal A_\infty$ the full fluctuation row;
the subscript $s$ includes the signed action on bonds. Linearity of
$C_0$ in its fine input and commutation of the replica with the signed
action give the displayed equation. In particular the cancellation
holds for each row, not just after summing the energy.

We will also need a volume-independent bound
\begin{equation}\label{eq:mean-response-cap}
 \|v_F(b)\|_\infty\le C_{\rm resp}\|b\|_\infty.
\end{equation}
Here is the quantitative tree estimate behind this bound. Put
$r_0=(J+3)L$ and let the installed cross kernel have weighted row
price $C/L$ and width $\eta\ge0$. At fine sites within $tL$ of a
fixed origin, moving the exponential weight costs at most $e^{\eta t}$.
Consequently its action on a scalar coarse field bounded by that
weight is at most
\[
 A/L,\qquad
 A=e^{\eta(12(J+3)+5)}e^\eta\,4\max(C,1).
\]
A tree hike has length at most $4L$. Each coordinate route in the
root patch has length at most $8r_0$; the root correction is the
difference of two averages of such path integrals. It is therefore
bounded by $16r_0 A/L$. Combining the hike and root correction gives
$C_1=(4+16(J+3))A$ for the straightening potential. Taking its bond
gradient costs $2C_1$, so the extracted response costs at most
$A+2C_1$. Subtracting the replica costs another $e^\eta$ in the
weighted estimate. For a field of sup norm one all its colour
coordinates satisfy the weighted premise, since the weight is at
least one. Summing the absolute colour coordinates bounds the
physical colour norm. Thus one may take
\[
 C_{\rm resp}=\max\{1,q\max(0,A+2C_1+e^\eta)\},
\]
where $q$ is the colour dimension. Homogeneity proves
\eqref{eq:mean-response-cap} for arbitrary $b$. All tree and path
lengths were bounded before periodization, so the estimate is uniform
in the period and in the forest.

\subsection{A common logarithmic comparison field}
Retain the upper-exponent convention of
Section~\ref{sec:packet-estimates}. Away from the seed set the
coarse field $V$ is clean: its plaquette chords are bounded by the
small threshold. If the ball of radius $d+1$ about $o$ is clean,
$d\le32B$, the axial construction gives a gauge with bond chords,
on the radius-$d$ ball, at most
\[
 16\{2(32B+1)+1\}^5N_p\,2\tau_c.
\]
Here $\tau_c$ is the coarse small threshold and has upper exponent
$-1/2$. Apply the local logarithm there and extend the logarithmic
field by zero elsewhere. We obtain a single global field $b$ and
a coarse gauge $g$ such that
\begin{equation}\label{eq:mean-common-log}
 g\cdot V=\exp b\ \hbox{on the radius-$d$ ball},\qquad
 \|b\|_\infty\le w(x),
\end{equation}
where $w\ge0$ has upper exponent $-1/2$. The period and chart-radius
conditions hold at one threshold depending only on the fixed step
parameters. The zero extension is used for the full linear kernels;
no smallness of $V$ outside the ball is asserted.

Consider a row whose anchor is at distance greater than $20B$ from
all seeds, and suppose \eqref{eq:mean-common-log} holds on its
radius-$2B$ ball. In the original gauge use the candidate
\[
 v=\operatorname{Ad}(g^{-1}\circ\operatorname{parent})v_F(b).
\]
It vanishes on tree bonds and has norm at most
$W=C_{\rm resp}w$. Within the stencil of this remote row all non-tree
bonds are active. Thus the sector's affine signal agrees there with
$v$, including its tree zeros. Gauge covariance transfers its actual
affine row to the same $b$ and $v_F(b)$.

Let $C(V)$ be the fine-curvature/predictor pair at the coarse replica,
$G(V)$ the guide, and $\mathcal A_R^V$ the transported cut fluctuation
row. Their residual is
$C(V)-G(V)+\mathcal A_R^Vv_F(b)$ in this common gauge.
Subtract \eqref{eq:mean-stationary-zero}. The four errors are
\[
 \begin{split}
 &C(V)-C_0(\overline b,b),\qquad -(G(V)-G_0(b)),\\
 &(\mathcal A_R^V-\mathcal A_R)v_F(b),\qquad
 (\mathcal A_R-\mathcal A_\infty)v_F(b).
 \end{split}
\]
The first two have upper exponent $-1$, by the local Taylor and
guide comparisons of Sections~\ref{sec:seed-unhit-comparison}
and~\ref{sec:packet-annulus}: their bounds are fixed multiples of
$w\{(1+R)w+e^{-\sigma R}\}$. The clean condition makes the guide's
clipping inactive on precisely the read stencil; it is removed
before the comparison with the globally extended logarithm.
For the third error, all coarse transport bonds have chord at most
$Cw$ on the common hull. The adjoint estimate and the absolute path
price give a bound $C(B+1)wW$. The fourth is at most
$CW e^{-\sigma R}$. The cutoff $R=\lfloor B/100000\rfloor$ has
an exponentially small tail of every prescribed negative upper
exponent. We have therefore proved, for one $a(x)\ge0$ of upper
exponent $-1$,
\begin{equation}\label{eq:mean-candidate-residual}
 \|\mathcal R_j(v)\|\le a(x)
\end{equation}
on every row whose comparison hull is included in the common clean
ball. Here $\mathcal R_j$ is the actual unscaled affine row, with
its original orientation and tags. One field $b$ is used for all
rows in the ball, so these estimates can be inserted into a matrix
identity without changing the candidate from row to row.

\subsection{Comparing the candidate with the chosen local mean}
Use counting coordinates and the matrices of
Sections~\ref{sec:sector-gram} and~\ref{sec:local-covariance}.
Write $M$ for the selected row matrix, $A=M^*M+\Pi_{\rm inactive}$,
$\ell=M^*r(0)$ for the linear source, and
\[
 r(v)=Mv+r(0),\qquad P=\hbox{the padded principal inverse},
 \qquad m=-P^*\ell,\quad E=AP-1.
\]
The row $r_j(v)$ is $s\mathcal R_j(v)$ when selected, and zero
otherwise; $s>0$ is the fixed signed-channel normalization. The
chosen physical mean $\mu$ is obtained from $m$ by reassembling its
colour coordinates. We have the exact identity
\begin{equation}\label{eq:mean-error-identity}
 m-v=E^*v-P^*(Av+\ell),\qquad
 Av+\ell=M^*r(v)+\Pi_{\rm inactive}v.
\end{equation}
For the first equality use $E^*=P^*A-1$ and cancel $P^*Av$.
The second follows directly from the definition of $A$ and $\ell$.

Let $S$ bound both unweighted row and column sums of $P$, let
$\varepsilon_B$ bound those of $E$, and let $Y$ bound those of $M$.
These constants are supplied by the weighted estimates already proved;
$\varepsilon_B$ has every negative upper exponent. If the normal
residual $Av+\ell$ is bounded by $b_0$ on the principal window
attached to coordinate $i$, \eqref{eq:mean-error-identity} gives
\begin{equation}\label{eq:mean-point}
 |m_i-v_i|\le S b_0+\varepsilon_B\|v\|_\infty.
\end{equation}
Indeed the $i$th row of $P^*$ is supported on that very window,
and the column bounds give both estimates. This uses $P^*$ in the
mean; replacing it by the symmetrized covariance would change the
identity.

Fix a target row at anchor $o$ farther than $30B$ from every seed.
Choose \eqref{eq:mean-common-log} on the radius-$8B$ ball. For each
coordinate within $B$ of $o$, its principal window lies within $3B$
of $o$, since the window radius is at most $B+2$. At such window
coordinates the inactive term vanishes: a non-tree bond is active
there, and a tree bond has candidate value zero. A nonzero entry
of $M^*$ reads a row within $B/4$, so all its affine residuals obey
\eqref{eq:mean-candidate-residual}, using subballs of the same
radius-$8B$ comparison. Hence $b_0=Ys a$, and
\[
 |m_i-v_i|\le U(x):=SYs\,a(x)+\varepsilon_BW(x)
 \quad\hbox{when }d(o,i)\le B.
\]
The function $U$ has upper exponent $-1$.

The target row is selected and reads only coordinates within $B/4$.
Its difference on $\mu$ and $v$ is linear, so the row bound on $M$
and the passage from colour coordinates to the colour norm imply
\[
 \|\mathcal R_o(\mu)-\mathcal R_o(v)\|
       \le q s^{-1}YU.
\]
Combining this with \eqref{eq:mean-candidate-residual} proves
\begin{equation}\label{eq:mean-remote-final}
 \|\mathcal R_o(\mu)\|\le a(x)+q s^{-1}YU(x),
 \qquad d(o,\text{seeds})>30B.
\end{equation}
The right-hand side has upper exponent $-1$, uniformly in the
period, forest, tags and all sectors satisfying the fixed construction
conditions. This is the required extra order of smallness of the
actual affine mean. The argument uses the flat stationary response
and its diagonal identity; it makes no assertion about derivatives
of the mean or a nonlinear global stationary solution.
